\subsection{Zadanie nr 1}

\subsubsection*{Treść}

W chwili $t = \SI{1.5}{s}$:
\begin{enumerate}
    \item Narysuj układ wektorów prędkości i przyspieszenia,
    \item Wyznacz składową normalną i styczną przyspieszenia,
    \item Określ krzywiznę toru.
\end{enumerate}

Ruch punktu dany jest równaniem parametrycznym:
\begin{align*}
    \Vec{r} &= [ x(t), y(t) ] \\
    x(t) &= \sin(\pi t) \\
    y(t) &= t^2
\end{align*}

\subsubsection*{Rozwiązanie}

Wyznaczamy wektory prędkości i przyspieszenia przez obliczenie pochodnych po czasie wektora położenia:

\begin{align*}
    x(t) &= \sin(\pi t) \\
    v_x(t) &= \dot x(t) = \pi \cos(\pi t) \\
    a_x(t) &= \dot v_x(t) = - \pi^2 \sin(\pi t) = -9.8696 \sin(\i t)
\end{align*}

\begin{align*}
    y(t) &= t^2 \\
    v_y(t) &= \dot y(t) = 2t \\
    a_y(t) &= \dot v_y(t) = 2
\end{align*}

Oznacza to, że w chwili $t = \SI{1.5}{s}$ wektory prędkości i przyspieszenia mają postać:

\begin{align*}
    \vec v(\SI{1.5}{s}) &= \left[ \pi \cos\left( \dfrac{3}{2} \pi \right), 2 \cdot 1.5 \right] = [0, 3] \ \si{m/s} \\
    \vec a(\SI{1.5}{s}) &= \left[ -\pi^2 \sin\left( \dfrac{3}{2} \pi \right), 2 \right] = [ -9.8696, 2] \ \si{m/s^2}
\end{align*}

